\section{The Postcard, Audited}

This section does something the rest of the article has not done: it checks the popular picture against the text, item by item, and reports the results of bounded searches. The method is crude and its limits are stated. What comes out of it is not a debunking. It is a redistribution---several of the postcard's furnishings turn out to be in Scripture after all, in odder places than expected, and the one thing the postcard never shows turns out to be the only thing the text puts at the centre.

\subsection{Method and its limits}

Every search reported here was run over the source library's tracked per-verse text of the Douay--Rheims in the Challoner electronic edition, on 2026-07-26. The searches are literal English string matches within single verses of that edition. They cannot exclude synonyms, paraphrase, a different translation's vocabulary, an inflection not searched, or a sense expressed without the word. Where a search returns nothing, the honest form of the result is ``this word does not occur in this edition in this range,'' and nothing stronger.\endnote{The tracked artifacts are the per-book, one-verse-per-line normalised text files registered under the Challoner Project Gutenberg edition of the Douay--Rheims; the New Testament range is the twenty-seven files numbered 47 through 73. This edition preserves the Vulgate chapter-and-verse numbering, which differs from a modern Bible's at several places relevant here---most importantly 1 Thess 4, where the ``caught up in the clouds'' verse is numbered 16 rather than 17. The edition also differs from the 1899 American edition in 1859 verses, collated in a table registered with it; none of the divergences affects any claim in this section.}

\subsection{What is in the text}

Taking the Apocalypse's final vision first, since it is the fullest description Scripture gives: a new heaven and a new earth, with the sea gone; a holy city, new Jerusalem, coming down out of heaven, prepared as a bride adorned for her husband; the tabernacle of God with men; every tear wiped away, and death, mourning, crying and sorrow ended, ``for the former things are passed away.''\endnote{Apoc 21:1--4, tracked verse text, read 2026-07-26.} Then the measurements: a wall great and high with twelve gates and twelve angels in them, the names of the twelve tribes on the gates and the twelve apostles on the foundations; the city four-square, twelve thousand furlongs, length and height and breadth equal; the wall of jasper, the city of pure gold like clear glass, the foundations adorned with twelve named gems, the twelve gates twelve pearls, the street pure gold ``as it were, transparent glass.''\endnote{Apoc 21:12--21, tracked verse text, read 2026-07-26.}

Then the four negations that carry the theology: ``And I saw no temple therein. For the Lord God Almighty is the temple thereof, and the Lamb.'' ``And the city hath no need of the sun, nor of the moon, to shine in it.'' ``And the gates thereof shall not be shut by day: for there shall be no night there.'' And in the next chapter, ``And there shall be no curse any more.''\endnote{Apoc 21:22, 23, 25; 22:3, tracked verse text, read 2026-07-26. The Douay of 21:23 continues: ``For the glory of God hath enlightened it: and the Lamb is the lamp thereof''---the verse Aquinas cites for the light of glory at ST I, q.~12, a.~5.} Then the river of the water of life clear as crystal, the tree of life bearing twelve fruits and yielding every month, its leaves ``for the healing of the nations.''\endnote{Apoc 22:1--2, tracked verse text, read 2026-07-26. The 1899 American edition reads the last clause slightly differently (``and the leaves of the tree were for the healing of the nations''), one of the collated divergences.} And then the sentence everything has been building to: ``And they shall see his face: and his name shall be on their foreheads.''\endnote{Apoc 22:4, tracked verse text, read 2026-07-26.}

Elsewhere: white robes and palms in the hands of a multitude no man could number, from every nation and tribe and people and tongue.\endnote{Apoc 7:9, tracked verse text, read 2026-07-26.} Crowns, repeatedly---a crown of justice, a crown of life, a never-fading crown of glory.\endnote{2 Tim 4:8; James 1:12; 1 Pet 5:4; Apoc 2:10; tracked verse text, read 2026-07-26.} Brightness: ``Then shall the just shine as the sun, in the kingdom of their Father''; ``they that are learned, shall shine as the brightness of the firmament: and they that instruct many to justice, as stars for all eternity.''\endnote{Matt 13:43; Dan 12:3; tracked verse text, read 2026-07-26.} A new heavens and a new earth ``in which justice dwelleth,'' promised in both Testaments.\endnote{2 Pet 3:13; Is 65:17; tracked verse text, read 2026-07-26.}

And harps, which the postcard did not invent. In the Apocalypse harps appear three times: the four living creatures and the twenty-four ancients hold them before the Lamb; a voice from heaven sounds ``as the voice of harpers, harping on their harps''; and those who have overcome the beast stand on a sea of glass mingled with fire, ``having the harps of God.''\endnote{Apoc 5:8; 14:2; 15:2; tracked verse text, read 2026-07-26. The literal search for the string ``harp'' over the tracked Apocalypse artifact was run with the source library's search tool on 2026-07-26; it returns these three plus a set of false positives from the substring in ``sharp,'' which is a fair illustration of what a literal line search can and cannot do.} The redeemed do get harps. What they do not get is one apiece on a cloud.

\subsection{What is not in the text}

Four items in the standard picture were tested and are not there in the form the picture gives them.

\emph{Wings.} A literal search for ``wings'' across the twenty-seven New Testament books of this edition returns five verses. Two are the Lord's lament over Jerusalem, where the wings are a hen's or a bird's. One gives six wings to the four living creatures around the throne. One gives wings to the locusts of the fifth trumpet. One gives ``two wings of a great eagle'' to the woman of chapter twelve, so that she may fly into the desert. In this edition's New Testament, no human being, living or risen or blessed, is given wings.\endnote{Matt 23:37; Luke 13:34; Apoc 4:8; 9:9; 12:14; searched and read 2026-07-26 over the tracked per-verse artifacts for the twenty-seven New Testament books. The result is bounded as stated in the method note above: it is the absence of an English word in one edition's New Testament, not a demonstration that no scriptural text anywhere assigns wings to the blessed. The winged human figure is an artistic convention of considerable antiquity and no scriptural warrant that this search located; its history is outside this article's scope.}

\emph{Clouds as a dwelling.} Clouds are everywhere in the New Testament, but almost always as the vehicle or veil of a divine coming: the cloud at the Transfiguration, the cloud that received Christ at the Ascension, the Son of man coming in the clouds. Twice human beings are associated with a cloud in the way the picture suggests, and both are movements rather than residences: the two witnesses of chapter eleven go up to heaven in a cloud, and at the parousia ``we who are alive, who are left, shall be taken up together with them in the clouds to meet Christ, into the air: and so shall we be always with the Lord.''\endnote{Apoc 11:12; 1 Thess 4:16 in the Vulgate numbering of this edition (4:17 in most modern Bibles); tracked verse text, read 2026-07-26. Acts 1:9, Matt 17:5, and Matt 24:30 are the theophanic uses referred to. The clause that matters in the Thessalonians text is the last one: the destination is not the cloud but ``always with the Lord.''} The cloud is transport. Nobody lives on it.

\emph{A gate-keeper.} The gates of the city have twelve angels stationed in them and the names of the twelve tribes written on them; the foundations carry the names of the apostles. But the gates are never shut, and nothing in the vision puts anyone at them with a register. The keys of the kingdom are given to Peter on earth, in a text about binding and loosing.\endnote{Apoc 21:12, 14, 25; Matt 16:19; tracked verse text, read 2026-07-26.} The comic scene of the interview at the gate is a joke about the last judgment relocated to the door of heaven, and the last judgment, in the texts, happens somewhere else and to everyone at once.

\emph{Becoming angels.} Both Synoptic forms of the Sadducees' question say something precise and less than the picture claims. Matthew: ``in the resurrection they shall neither marry nor be married, but shall be \emph{as} the angels of God in heaven.'' Luke: ``Neither can they die any more for they are \emph{equal to} the angels and are the children of God, being the children of the resurrection.''\endnote{Matt 22:30; Luke 20:36; tracked verse text, read 2026-07-26; emphasis added. Augustine makes the point that the comparison is limited: ``They shall be equal to the angels in immortality and happiness, not in flesh, nor in resurrection, which the angels did not need, because they could not die'' (\emph{De civitate Dei} XXII.17, Dods translation, read in the tracked transcription 2026-07-26).} Both are comparisons, and both name their term: not marrying, not dying. Neither says the risen become angels, and the whole doctrine of the resurrection of the body says they do not.

One further item, of a different kind: the string ``halo'' does not occur in this edition of Scripture at all except inside proper names.\endnote{Searched 2026-07-26 over the tracked per-verse artifacts; the seven matches are all substrings of proper nouns (Mahalon, Chelion's neighbours in Ruth, Hethalon in Ezechiel, and the like). The nimbus is an iconographic convention, and a good one; it is not a report.}

\subsection{What kind of language this is}

It would be a poor conclusion to say that Scripture's heaven is imagery and therefore tells us nothing. The Catechism draws the opposite conclusion from the same fact: ``This mystery of blessed communion with God and all who are in Christ is beyond all understanding and description. Scripture speaks of it in images: life, light, peace, wedding feast, wine of the kingdom, the Father's house, the heavenly Jerusalem, paradise.''\endnote{CCC 1027, official Vatican English web text, read 2026-07-26. The paragraph closes by quoting 1 Cor 2:9.} The images are revealed. They are not a concession to weak minds that a stronger mind could dispense with; they are the form in which the content is given, and the content is genuinely in them.

The tradition has always known this and said it when pressed. Augustine, in the middle of his most speculative chapter, has to stop and explain that ``by `the face' of God we are to understand His manifestation, and not a part of the body similar to that which in our bodies we call by that name.''\endnote{\emph{De civitate Dei} XXII.29, Dods translation, read 2026-07-26 in the tracked transcription.} Aquinas, denying that God can be seen in an imaginative vision, notes that ``in the divine Scripture divine things are metaphorically described by means of sensible things.''\endnote{ST I, q.~12, a.~3, reply to objection 3, read 2026-07-26.} And Paul, who says he was caught up to the third heaven and into paradise, reports that he ``heard secret words which it is not granted to man to utter''---the one first-hand account in the canon, and it declines to give an account.\endnote{2 Cor 12:2--4, tracked verse text, read 2026-07-26. Compare 1 Cor 2:9: ``That eye hath not seen, nor ear heard: neither hath it entered into the heart of man, what things God hath prepared for them that love him,'' itself a citation of Is 64:4.}

Read as symbols, the images give more than the postcard, not less. No temple: the distinction between sacred and secular time, sacred and secular space, is abolished---not because worship stops but because nothing is outside it. No sun and no night: the city's light is not something it receives from a body it orbits. Gates never shut: a walled city with permanently open gates is a city with nothing left to fear, which is what a fortified gate is for. Every tear wiped away, by a hand, personally. A tree whose leaves are for healing, in a place where nothing is sick. And names---on the gates, on the foundations, on the foreheads of the blessed. Identity, twelve times over, written into the architecture.

\subsection{The one thing the postcard leaves out}

Line up the images and a structure appears. Almost all of them are negations of the conditions of this life: no sea, no death, no mourning, no crying, no sorrow, no temple, no sun, no night, no curse, no shut gate. The Apocalypse describes heaven mainly by subtraction, which is exactly what one would expect of a text describing something for which its readers have no positive concept.

Against that there stands one positive clause, and it is the last thing said before the city's lamps are described: \emph{and they shall see his face}.

That is the sentence the definitions of 1336 and 1439 were written to protect, and it is the sentence the postcard has no way to draw. A picture can show clouds, harps, gates, and reunions; it cannot show the vision of the divine essence, because there is nothing to show---the seeing is an act of a mind and its object has no shape. So the picture fills the space with the furniture, and the furniture, being visible, becomes what heaven is about.

That is the whole mechanism of the drift, and it explains why the drift runs the way it does. It is not that Christians became more worldly or less devout. It is that the central content of the doctrine is the one item that cannot be illustrated, and illustration is how a picture works. The corrective is not to burn the postcards. It is to know what the postcards are for, and to know that the small clause they omit is the only part the Church has ever bound anyone to believe.
